\section{Характеристики эффективности системы M|M|1}

Ранее было установлено, что стационарное распределение вероятностей состояний классической одноканальной системы M|M|1 существует при условии $\rho = \frac{\lambda}{\mu} < 1$ и задается геометрическим законом: $p_n = (1 - \rho)\rho^n$. Опираясь на этот результат, определим ключевые макроскопические метрики эффективности функционирования системы.

\begin{statement}[Стационарные характеристики системы M|M|1]
Для одноканальной системы массового обслуживания с пуассоновским входным потоком интенсивности $\lambda$ и экспоненциальным временем обслуживания интенсивности $\mu$ при $\lambda < \mu$ справедливы следующие соотношения:
\begin{itemize}
    \item Среднее число заявок в системе: $L = \frac{\rho}{1 - \rho}$.
    \item Среднее число заявок в очереди: $L_q = \frac{\rho^2}{1 - \rho}$.
    \item Среднее время пребывания заявки в системе: $W = \frac{1}{\mu - \lambda}$.
    \item Среднее время ожидания заявки в очереди: $W_q = \frac{\rho}{\mu - \lambda}$.
\end{itemize}
\end{statement}

\begin{sproof}
\textbf{1. Среднее число заявок в системе ($L$).} По определению математического ожидания дискретной случайной величины $N$ имеем:
$$
L = \E[N] = \sum_{n=0}^{\infty} n \cdot p_n = (1 - \rho) \sum_{n=0}^{\infty} n \rho^n.
$$
Для вычисления суммы ряда воспользуемся методом дифференцирования по параметру $\rho$. Вынесем $\rho$ за знак суммы:
$$
\sum_{n=0}^{\infty} n \rho^n = \rho \sum_{n=1}^{\infty} n \rho^{n-1} = \rho \frac{\partial}{\partial \rho} \left( \sum_{n=1}^{\infty} \rho^n \right).
$$
Поскольку сумма бесконечной геометрической прогрессии равна $\frac{\rho}{1-\rho}$, находим производную:
$$
\frac{\partial}{\partial \rho} \left( \frac{\rho}{1-\rho} \right) = \frac{1 \cdot (1-\rho) - \rho \cdot (-1)}{(1-\rho)^2} = \frac{1}{(1-\rho)^2}.
$$
Подставляя полученный результат в исходное выражение для $L$, получаем:
$$
L = (1 - \rho) \cdot \rho \cdot \frac{1}{(1-\rho)^2} = \frac{\rho}{1-\rho} = \frac{\lambda}{\mu - \lambda}.
$$

\textbf{2. Среднее число заявок в очереди ($L_q$).} Длина очереди $N_q$ связана с общим числом заявок соотношением $N_q = \max(0, N - 1)$. Используя линейность математического ожидания:
$$
L_q = \E[N_q] = \sum_{n=1}^{\infty} (n - 1) p_n = \sum_{n=1}^{\infty} n p_n - \sum_{n=1}^{\infty} p_n = \E[N] - (1 - p_0).
$$
Так как вероятность простоя системы $p_0 = 1 - \rho$, то $1 - p_0 = \rho$. Отсюда:
$$
L_q = L - \rho = \frac{\rho}{1-\rho} - \rho = \frac{\rho - \rho(1-\rho)}{1-\rho} = \frac{\rho^2}{1-\rho} = \frac{\lambda^2}{\mu(\mu - \lambda)}.
$$

\textbf{3. Временные характеристики ($W$ и $W_q$).} Применим формулу Литтла ($L = \lambda W$):
$$
W = \frac{L}{\lambda} = \frac{\lambda}{\lambda(\mu - \lambda)} = \frac{1}{\mu - \lambda}.
$$
Аналогично для очереди ($L_q = \lambda W_q$):
$$
W_q = \frac{L_q}{\lambda} = \frac{\lambda^2}{\lambda \mu (\mu - \lambda)} = \frac{\lambda}{\mu(\mu - \lambda)}.
$$
\end{sproof}

Отдельный интерес представляет исследование длины очереди при условии, что она физически существует (поступающая заявка застает прибор занятым).

\begin{conseq}
Среднее число заявок в очереди при условии, что очередь не пуста, равно:
$$
L_q' = \frac{1}{1-\rho}.
$$
\end{conseq}

\begin{sproof}
Введем условное математическое ожидание $L_q' = \E[N_q \mid N_q > 0]$. Очередь не пуста, когда в системе находится $n \ge 2$ заявок. Вероятность этого события:
$$
\P(N \ge 2) = 1 - p_0 - p_1 = 1 - (1-\rho) - (1-\rho)\rho = \rho^2.
$$
По определению условного математического ожидания:
$$
L_q' = \sum_{n=2}^{\infty} (n-1) \frac{p_n}{\P(N \ge 2)} = \frac{1}{\rho^2} \sum_{n=2}^{\infty} (n-1)p_n.
$$
Ряд в числителе совпадает с безусловным ожиданием $L_q$ (так как при $n=1$ слагаемое зануляется):
$$
L_q' = \frac{L_q}{\rho^2} = \frac{1}{\rho^2} \cdot \frac{\rho^2}{1-\rho} = \frac{1}{1-\rho}.
$$
\end{sproof}

\begin{note}
Зависимость метрик $L, L_q, W, W_q$ от коэффициента загрузки $\rho$ имеет гиперболический характер с вертикальной асимптотой при $\rho = 1$. Это означает неограниченный рост очередей и задержек при приближении загрузки к 100\%.
\end{note}

\begin{center}
\begin{minipage}{0.48\textwidth}
    \centering
    \begin{tikzpicture}
    \begin{axis}[
        width=\linewidth,
        height=6cm,
        xlabel={Загрузка $\rho$},
        ylabel={Среднее число заявок},
        xmin=0, xmax=1,
        ymin=0, ymax=10,
        grid=major,
        legend pos=north west,
        legend style={font=\footnotesize},
        tick label style={font=\footnotesize},
        label style={font=\small}
    ]
        % L = rho / (1-rho)
        \addplot[DeepSkyBlue4, thick, domain=0:0.9] {x/(1-x)};
        \addlegendentry{$L$}
        
        % Lq = rho^2 / (1-rho)
        \addplot[Red3, thick, dashed, domain=0:0.9] {x^2/(1-x)};
        \addlegendentry{$L_q$}
        
        \draw[dashed, gray] (axis cs:1,0) -- (axis cs:1,10);
    \end{axis}
    \end{tikzpicture}
    \captionof*{figure}{\small Число заявок в системе и очереди}
\end{minipage}
\hfill
\begin{minipage}{0.48\textwidth}
    \centering
    \begin{tikzpicture}
    \begin{axis}[
        width=\linewidth,
        height=6cm,
        xlabel={Загрузка $\rho$},
        ylabel={Норм. время $\mu W, \mu W_q$},
        xmin=0, xmax=1,
        ymin=0, ymax=10,
        grid=major,
        legend pos=north west,
        legend style={font=\footnotesize},
        tick label style={font=\footnotesize},
        label style={font=\small}
    ]
        % W*mu = 1 / (1-rho)
        \addplot[DeepSkyBlue4, thick, domain=0:0.9] {1/(1-x)};
        \addlegendentry{$W \cdot \mu$}
        
        % Wq*mu = rho / (1-rho)
        \addplot[Red3, thick, dashed, domain=0:0.9] {x/(1-x)};
        \addlegendentry{$W_q \cdot \mu$}
        
        \draw[dashed, gray] (axis cs:1,0) -- (axis cs:1,10);
    \end{axis}
    \end{tikzpicture}
    \captionof*{figure}{\small Временные характеристики}
\end{minipage}
\end{center}

\subsection*{Распределение времени ожидания и пребывания в системе M|M|1}

До сих пор макроскопические характеристики системы вычислялись без учета дисциплины обслуживания. Однако для анализа распределения времени ожидания конкретной заявки необходимо зафиксировать дисциплину очереди. Предположим, что заявки обслуживаются в порядке поступления (FCFS — First-Come, First-Served). 

Обозначим через $T_q$ случайную величину времени ожидания заявки в очереди, а через $T$ — полное время ее пребывания в системе. Соответствующие функции распределения задаются как $W_q(t) = \P(T_q \le t)$ и $W(t) = \P(T \le t)$.

\begin{statement}[Функции распределения времени задержки M|M|1]
Для системы M|M|1 с дисциплиной обслуживания FCFS функции распределения времени ожидания и времени пребывания имеют вид:
\begin{align*}
    W_q(t) &= \begin{cases}
        0, & t < 0, \\
        1 - \rho, & t = 0, \\
        1 - \rho e^{-(\mu - \lambda)t}, & t > 0;
    \end{cases} \\[2mm]
    W(t) &= \begin{cases}
        0, & t \le 0, \\
        1 - e^{-(\mu - \lambda)t}, & t > 0.
    \end{cases}
\end{align*}
\end{statement}

\begin{sproof}
Случайная величина $T_q$ имеет смешанное распределение. Пусть $N$ — число заявок в системе в момент поступления новой заявки. Применим формулу полной вероятности:
$$
W_q(t) = \P(T_q \le t) = \sum_{n=0}^{\infty} \P(T_q \le t \mid N = n) \P(N = n).
$$

В силу свойства PASTA вероятности состояний в момент прибытия совпадают со стационарными: $\P(N = n) = p_n$.

С вероятностью $p_0 = 1 - \rho$ поступающая заявка застает систему пустой ($N = 0$) и немедленно принимается к обслуживанию. В этом случае время ожидания равно нулю, и $\P(T_q \le t \mid N = 0) = 1$ для любого $t \ge 0$. Отсюда возникает дискретный скачок функции распределения в нуле: 
$$
W_q(0) = \P(T_q = 0) = p_0 = 1 - \rho.
$$

Пусть $t > 0$ и $n \ge 1$. Поскольку рассматривается FCFS, заявка перейдет к обслуживанию за время $\le t$, если суммарное время обслуживания всех $n$ заявок, находящихся перед ней, не превзойдет $t$. Так как времена обслуживания независимы и экспоненциально распределены, время суммарного обслуживания $n$ заявок подчиняется распределению Эрланга порядка $n$ с плотностью $f_n(x) = \frac{\mu^n}{(n-1)!} x^{n-1} e^{-\mu x}$. Таким образом:
$$
\P(T_q \le t \mid N = n) = \int_0^t f_n(x) dx.
$$

Подставляя полученные выражения в формулу полной вероятности, имеем:
$$
W_q(t) = 1 - \rho + \sum_{n=1}^{\infty} p_n \int_0^t f_n(x) dx = 1 - \rho + \sum_{n=1}^{\infty} (1 - \rho)\rho^n \int_0^t \frac{\mu^n x^{n-1}}{(n-1)!} e^{-\mu x} dx.
$$

Поменяем порядок суммирования и интегрирования, вынеся независимые от $n$ множители за знак суммы:
$$
W_q(t) = 1 - \rho + \int_0^t \mu (1 - \rho) \rho e^{-\mu x} \left( \sum_{n=1}^{\infty} \frac{(\mu \rho x)^{n-1}}{(n-1)!} \right) dx.
$$

Сделав замену индекса суммирования $k = n - 1$, замечаем, что ряд в скобках сводится к разложению экспоненты:
$$
\sum_{n=1}^{\infty} \frac{(\mu \rho x)^{n-1}}{(n-1)!} = \sum_{k=0}^{\infty} \frac{(\mu \rho x)^k}{k!} = e^{\mu \rho x} = e^{\lambda x}.
$$

Подставив свернутый ряд в интеграл, находим:
$$
W_q(t) = 1 - \rho + \rho \int_0^t \mu (1 - \rho) e^{-(\mu - \lambda)x} dx = 1 - \rho + \rho \left( 1 - e^{-(\mu - \lambda)t} \right) = 1 - \rho e^{-(\mu - \lambda)t}.
$$

Проверим согласованность результата, вычислив математическое ожидание времени ожидания через интеграл от функции надежности:
$$
W_q = \E[T_q] = \int_0^{\infty} (1 - W_q(t)) dt = \int_0^{\infty} \rho e^{-(\mu - \lambda)t} dt = \frac{\rho}{\mu - \lambda}.
$$
Полученное значение совпадает с формулой, выведенной ранее из закона Литтла.

Для функции распределения полного времени пребывания в системе $W(t)$ логика аналогична. Отличие состоит в том, что для ухода заявки из системы к моменту времени $t$ должно завершиться обслуживание $n+1$ заявок (включая саму прибывшую). Учитывая отсутствие дискретной компоненты в нуле ($W(0) = 0$), суммирование плотности Эрланга порядка $n+1$ с весами $p_n$ дает строго непрерывное экспоненциальное распределение:
$$
W(t) = 1 - e^{-(\mu - \lambda)t}.
$$
\end{sproof}

\begin{center}
    \pgfplotsset{compat=1.9}
    \begin{tikzpicture}
        \begin{axis}[
            width=0.85\textwidth, height=7cm,
            xlabel={Время $t$}, ylabel={$\P(\cdot \le t)$},
            xmin=0, xmax=5, ymin=0, ymax=1.1,
            grid=major,
            legend pos=south east,
            legend style={font=\small},
            title={Функции распределения (для $\mu=2, \lambda=1 \Rightarrow \rho=0.5$)}
        ]
            % W(t) - пребывание
            \addplot[Red3, very thick, domain=0:5, samples=100] expression {1 - exp(-(2-1)*x)};
            \addlegendentry{$W(t)$ (Пребывание)}
            
            % W_q(t) - ожидание
            \addplot[DeepSkyBlue4, very thick, domain=0:5, samples=100] expression {1 - 0.5*exp(-(2-1)*x)};
            \addlegendentry{$W_q(t)$ (Ожидание)}
            
            % Визуализация скачка функции в нуле
            \draw[DeepSkyBlue4, fill=white, thick] (axis cs:0, 0.5) circle (2pt);
            \draw[DeepSkyBlue4, fill=DeepSkyBlue4, thick] (axis cs:0, 0.5) circle (1.5pt);
            
        \end{axis}
    \end{tikzpicture}
    
    \vspace{2mm}
    \small\textit{Сравнение функций распределения времени ожидания и пребывания.}
\end{center}

\section{Многоканальная система обслуживания M|M|c}

Обобщим результаты на случай многоканальной системы, состоящей из $c$ идентичных параллельных серверов. На вход поступает пуассоновский поток заявок с интенсивностью $\lambda$. Времена обслуживания на серверах независимы и распределены экспоненциально с параметром $\mu$. 

Напомним ключевые характеристики нагрузки: 
\begin{itemize}
    \item $r = \frac{\lambda}{\mu}$ — приведенная интенсивность потока (нагрузка в Эрлангах);
    \item $\rho = \frac{r}{c} = \frac{\lambda}{c\mu}$ — коэффициент загрузки системы.
\end{itemize}

\begin{statement}[Стационарное распределение системы M|M|c]
Необходимым и достаточным условием существования стационарного режима является $\rho < 1$. Вероятности состояний $p_n$ задаются формулами:
$$
p_n = 
\begin{cases} 
    \dfrac{r^n}{n!} p_0, & 0 \le n < c, \\[3mm]
    \dfrac{r^n}{c! c^{n-c}} p_0, & n \ge c,
\end{cases}
$$
где вероятность простоя системы:
$$
p_0 = \left( \sum_{n=0}^{c-1} \frac{r^n}{n!} + \frac{r^c}{c!(1 - \rho)} \right)^{-1}.
$$
\end{statement}

\begin{sproof}
Используется общее решение для процессов рождения и гибели. Для $n < c$ интенсивность обслуживания равна $n\mu$, для $n \ge c$ она стабилизируется на уровне $c\mu$. Вывод формул проводится прямой подстановкой интенсивностей в произведение $\prod \frac{\lambda_{i-1}}{\mu_i}$ и суммированием геометрической прогрессии для нахождения константы нормировки $p_0$, аналогично тому, как это было сделано ранее для M|M|1.
\end{sproof}

\begin{statement}[Стационарные характеристики системы M|M|c]
Для многоканальной системы массового обслуживания, состоящей из $c$ приборов, с пуассоновским входным потоком интенсивности $\lambda$ и экспоненциальным временем обслуживания интенсивности $\mu$ при коэффициенте загрузки $\rho = \frac{\lambda}{c\mu} < 1$ справедливы следующие соотношения:
\begin{itemize}
    \item Среднее число заявок в очереди: $\displaystyle L_q = \frac{p_0 r^c \rho}{c! (1-\rho)^2}$.
    \item Среднее число заявок в системе: $\displaystyle L  = L_q + r = \frac{p_0 r^c \rho}{c! (1-\rho)^2} + r$.
    \item Среднее время ожидания заявки в очереди: $\displaystyle W_q = \frac{L_q}{\lambda} = \frac{p_0 r^c}{c \cdot c! \mu (1-\rho)^2}$.
    \item Среднее время пребывания заявки в системе: $\displaystyle W = W_q + \frac{1}{\mu} = \frac{p_0 r^c}{c \cdot c! \mu (1-\rho)^2} + \frac{1}{\mu}$.
\end{itemize}
где $r = \frac{\lambda}{\mu}$ — приведенная интенсивность потока, а $p_0$ — вероятность простоя системы.
\end{statement}

\begin{sproof}
\textbf{1. Среднее число заявок в очереди ($L_q$).} По определению математического ожидания дискретной случайной величины длины очереди $N_q = \max(0, N - c)$ имеем:
$$
L_q = \E[N_q] = \sum_{n=c}^{\infty} (n - c) p_n = \sum_{n=c+1}^{\infty} (n - c) p_n.
$$
Подставим выражение для вероятностей состояний $p_n$ при $n \ge c$, заданное формулой $p_n = \frac{r^n}{c! c^{n-c}} p_0$, и введем новый индекс суммирования $m = n - c$:
$$
L_q = \sum_{m=1}^{\infty} m p_{c+m} = \sum_{m=1}^{\infty} m \frac{r^{c+m}}{c! c^m} p_0.
$$
Вынесем независимые от $m$ множители за знак суммы и сгруппируем члены, зависящие от $m$, используя соотношение $\rho = \frac{r}{c}$:
$$
L_q = \frac{p_0 r^c}{c!} \sum_{m=1}^{\infty} m \left(\frac{r}{c}\right)^m = \frac{p_0 r^c}{c!} \sum_{m=1}^{\infty} m \rho^m.
$$
Ряд $\dsum_{m=1}^{\infty} m \rho^m$ равен $\frac{\rho}{(1-\rho)^2}$. Подставляя это значение, получаем окончательное выражение для длины очереди:
$$
L_q = \frac{p_0 r^c \rho}{c! (1-\rho)^2}.
$$

\textbf{2. Временные характеристики ($W$ и $W_q$).} По закону Литтла, примененному непосредственно к подсистеме очереди, среднее время ожидания равно:
$$
W_q = \frac{L_q}{\lambda} = \frac{p_0 r^c \rho}{\lambda c! (1-\rho)^2}.
$$
Используя определение коэффициента загрузки $\rho = \frac{\lambda}{c\mu}$, заменяем отношение $\frac{\rho}{\lambda}$ на $\frac{1}{c\mu}$:
$$
W_q = \frac{p_0 r^c}{c \cdot c! \mu (1-\rho)^2}.
$$
Полное время пребывания заявки в системе $T$ складывается из времени ожидания в очереди $T_q$ и времени непосредственного обслуживания. Поскольку время обслуживания экспоненциально с параметром $\mu$, его математическое ожидание равно $\frac{1}{\mu}$. Из линейности математического ожидания следует:
$$
W = \E[T] = \E[T_q] + \E[\text{время обслуживания}] = W_q + \frac{1}{\mu} = \frac{p_0 r^c}{c \cdot c! \mu (1-\rho)^2} + \frac{1}{\mu}.
$$

\textbf{3. Среднее число заявок в системе ($L$).} Применим универсальный закон Литтла к системе в целом ($L = \lambda W$) и подставим найденное выражение для $W$:
$$
L = \lambda W = \lambda \left( W_q + \frac{1}{\mu} \right) = \lambda W_q + \frac{\lambda}{\mu}.
$$
Учитывая, что по закону Литтла для очереди $\lambda W_q = L_q$, а по определению $\frac{\lambda}{\mu} = r$, приходим к искомому результату:
$$
L = L_q + r = \frac{p_0 r^c \rho}{c! (1-\rho)^2} + r.
$$
\end{sproof}

\begin{example}{Аналитический расчет характеристик для малого числа каналов}
Продемонстрируем зависимость макроскопических метрик от $c$ в явном виде, выразив вероятность простоя $p_0$ и длину очереди $L_q$ исключительно через коэффициент загрузки $\rho$. Напомним, что приведенная нагрузка связана с загрузкой соотношением $r = c\rho$.

\textbf{Случай $c=1$ (одноканальная СМО):}
Вероятность простоя вычисляется сворачиванием геометрической прогрессии:
$$
p_0 = \left( 1 + \frac{\rho}{1-\rho} \right)^{-1} = (1 - \rho).
$$
Подставляя $p_0$ в общую формулу длины очереди, получаем классический результат:
$$
L_q^{(1)} = \frac{(1-\rho) \rho \cdot \rho}{1! (1-\rho)^2} = \frac{\rho^2}{1-\rho}.
$$

\textbf{Случай $c=2$ (двухканальная СМО):}
Ряд для нормировочной константы $p_0$ обрывается на втором слагаемом перед переходом к бесконечной сумме. Подставим $r = 2\rho$:
$$
p_0 = \left( 1 + \frac{2\rho}{1!} + \frac{(2\rho)^2}{2!(1-\rho)} \right)^{-1} = \left( 1 + 2\rho + \frac{2\rho^2}{1-\rho} \right)^{-1}.
$$
Приведем выражение в скобках к общему знаменателю:
$$
1 + 2\rho + \frac{2\rho^2}{1-\rho} = \frac{1 - \rho + 2\rho - 2\rho^2 + 2\rho^2}{1-\rho} = \frac{1 + \rho}{1-\rho} \implies p_0 = \frac{1-\rho}{1+\rho}.
$$
Подставляя найденное $p_0$ в формулу для очереди, находим:
$$
L_q^{(2)} = \frac{\frac{1-\rho}{1+\rho} (2\rho)^2 \rho}{2! (1-\rho)^2} = \frac{4\rho^3 (1-\rho)}{2(1+\rho)(1-\rho)^2} = \frac{2\rho^3}{1-\rho^2}.
$$

\textbf{Случай $c=4$ (четырехканальная СМО):}
Аналогичные алгебраические преобразования (суммирование четырех членов ряда и прогрессии для остатка) приводят к появлению полиномиальной дроби:
$$
p_0 = \frac{1-\rho}{1 + 3\rho + 4\rho^2 + \frac{8}{3}\rho^3}.
$$
Отсюда длина очереди выражается как:
$$
L_q^{(4)} = \frac{32 \rho^5}{3 (1-\rho) \left( 1 + 3\rho + 4\rho^2 + \frac{8}{3}\rho^3 \right)}.
$$

Именно эти нелинейные полиномиальные соотношения определяют поведение кривых на графиках эффективности. Можно заметить, что при фиксированной загрузке $\rho \to 1$ наличие дополнительных каналов не меняет порядка асимптоты $O((1-\rho)^{-1})$, но существенно уменьшает константу перед ней.
\end{example}

\begin{center}
\begin{minipage}{0.48\textwidth}
    \centering
    \pgfplotsset{compat=1.9}
    \begin{tikzpicture}
    \begin{axis}[
        width=\linewidth,
        height=7.5cm,
        xlabel={Загрузка $\rho$},
        ylabel={Среднее число заявок},
        xmin=0, xmax=1,
        ymin=0, ymax=10,
        grid=major,
        legend pos=north west,
        legend style={font=\scriptsize, cells={anchor=west}},
        legend columns=2,
        tick label style={font=\footnotesize},
        label style={font=\small}
    ]
        % c=1 (DeepSkyBlue4)
        \addplot[DeepSkyBlue4, thick, domain=0:0.9] {x/(1-x)};
        \addlegendentry{$L \ (c=1)$}
        \addplot[DeepSkyBlue4, thick, dashed, domain=0:0.9] {x^2/(1-x)};
        \addlegendentry{$L_q \ (c=1)$}
        
        % c=2 (Red3)
        \addplot[Red3, thick, domain=0:0.9] {2*x/(1-x^2)};
        \addlegendentry{$L \ (c=2)$}
        \addplot[Red3, thick, dashed, domain=0:0.9] {2*x^3/(1-x^2)};
        \addlegendentry{$L_q \ (c=2)$}
        
        % c=4 (SpringGreen4)
        \addplot[SpringGreen4, thick, domain=0:0.9] {(32/3*x^5)/((1-x)*(1+3*x+4*x^2+8/3*x^3)) + 4*x};
        \addlegendentry{$L \ (c=4)$}
        \addplot[SpringGreen4, thick, dashed, domain=0:0.9] {(32/3*x^5)/((1-x)*(1+3*x+4*x^2+8/3*x^3))};
        \addlegendentry{$L_q \ (c=4)$}
        
        \draw[dashed, gray] (axis cs:1,0) -- (axis cs:1,10);
    \end{axis}
    \end{tikzpicture}
    \captionof*{figure}{\small Число заявок в системе и очереди}
\end{minipage}
\hfill
\begin{minipage}{0.48\textwidth}
    \centering
    \pgfplotsset{compat=1.9}
    \begin{tikzpicture}
    \begin{axis}[
        width=\linewidth,
        height=7.5cm,
        xlabel={Загрузка $\rho$},
        ylabel={Норм. время $\mu W, \mu W_q$},
        xmin=0, xmax=1,
        ymin=0, ymax=10,
        grid=major,
        legend pos=north west,
        legend style={font=\scriptsize, cells={anchor=west}},
        legend columns=2,
        tick label style={font=\footnotesize},
        label style={font=\small}
    ]
        % c=1 (DeepSkyBlue4)
        \addplot[DeepSkyBlue4, thick, domain=0:0.9] {1/(1-x)};
        \addlegendentry{$W \ (c=1)$}
        \addplot[DeepSkyBlue4, thick, dashed, domain=0:0.9] {x/(1-x)};
        \addlegendentry{$W_q \ (c=1)$}
        
        % c=2 (Red3)
        \addplot[Red3, thick, domain=0:0.9] {1/(1-x^2)};
        \addlegendentry{$W \ (c=2)$}
        \addplot[Red3, thick, dashed, domain=0:0.9] {x^2/(1-x^2)};
        \addlegendentry{$W_q \ (c=2)$}
        
        % c=4 (SpringGreen4)
        \addplot[SpringGreen4, thick, domain=0:0.9] {(8/3*x^4)/((1-x)*(1+3*x+4*x^2+8/3*x^3)) + 1};
        \addlegendentry{$W \ (c=4)$}
        \addplot[SpringGreen4, thick, dashed, domain=0:0.9] {(8/3*x^4)/((1-x)*(1+3*x+4*x^2+8/3*x^3))};
        \addlegendentry{$W_q \ (c=4)$}
        
        \draw[dashed, gray] (axis cs:1,0) -- (axis cs:1,10);
    \end{axis}
    \end{tikzpicture}
    \captionof*{figure}{\small Время пребывания и ожидания}
\end{minipage}
\end{center}

\begin{definition}{Вероятность задержки}
Вероятность того, что поступающая заявка застанет все серверы занятыми и будет вынужден ожидать в очереди, обозначается $C(c, r)$.
\end{definition}

\begin{statement}[Формула Эрланга-C]
Вероятность задержки заявки вычисляется по формуле:
$$
C(c, r) \ \dot{=} \ \P(T_q > 0) = \frac{p_0 r^c}{c!(1-\rho)}.
$$
\end{statement}

\begin{sproof}
Заявка вынуждена ждать, если в момент ее прибытия в системе находится $n \ge c$ других заявок. Суммируя вероятности стационарного распределения:
$$
\P(T_q > 0) = \sum_{n=c}^{\infty} p_n = \sum_{n=c}^{\infty} \frac{r^n}{c! c^{n-c}} p_0 = \frac{p_0 r^c}{c!} \sum_{m=0}^{\infty} \rho^m = \frac{p_0 r^c}{c!(1-\rho)}.
$$
\end{sproof}

Это соотношение характеризует эффект масштаба: при фиксированном коэффициенте загрузки $\rho$ система с б\'{о}льшим числом серверов обеспечивает меньшую вероятность задержки, то есть обладает более высокой стабильностью под нагрузкой. Визуализация этого феномена представлена на графике ниже.

\begin{center}
    \definecolor{MyPurple}{RGB}{102, 51, 153}
    \pgfplotsset{compat=1.9}
    \begin{tikzpicture}
        \begin{axis}[
            width=0.85\textwidth,
            height=8cm,
            xlabel={Коэффициент загрузки $\rho$},
            ylabel={Вероятность задержки $C(c, c\rho)$},
            xmin=0, xmax=1,
            ymin=0, ymax=1.05,
            grid=major,
            legend pos=north west,
            legend columns=2,
            legend style={font=\footnotesize, fill=white, fill opacity=0.8, draw opacity=1},
            tick label style={font=\footnotesize},
            label style={font=\small},
            title={Вероятность ожидания (Формула Эрланга-C)}
        ]
            % c=1: C = x
            \addplot[Red3, very thick, domain=0:1, samples=50] {x};
            \addlegendentry{$c=1$}
            
            % c=2: C = 2x^2 / (1+x)
            \addplot[orange, very thick, domain=0:1, samples=50] {(2*x^2)/(1 + x)};
            \addlegendentry{$c=2$}
            
            % c=3
            \addplot[Tan1, very thick, domain=0:1, samples=50] 
            {(4.5*x^3) / ( (1-x)*(1 + 3*x + 4.5*x^2) + 4.5*x^3 )};
            \addlegendentry{$c=3$}
            
            % c=4
            \addplot[SpringGreen4, very thick, domain=0:1, samples=50] 
            {((32/3)*x^4) / ( (1-x)*(1 + 4*x + 8*x^2 + (32/3)*x^3) + (32/3)*x^4 )};
            \addlegendentry{$c=4$}
            
            % c=5
            \addplot[SeaGreen, very thick, domain=0:1, samples=50] 
            {((625/24)*x^5) / ( (1-x)*(1 + 5*x + 12.5*x^2 + (125/6)*x^3 + (625/24)*x^4) + (625/24)*x^5 )};
            \addlegendentry{$c=5$}
            
            % c=6 (Здесь 64.8 — конечное точное число, можно оставить десятичным)
            \addplot[DeepSkyBlue4, very thick, domain=0:1, samples=50] 
            {(64.8*x^6) / ( (1-x)*(1 + 6*x + 18*x^2 + 36*x^3 + 54*x^4 + 64.8*x^5) + 64.8*x^6 )};
            \addlegendentry{$c=6$}
            
            % c=7
            \addplot[RoyalBlue, very thick, domain=0:1, samples=50] 
            {((117649/720)*x^7) / ( (1-x)*(1 + 7*x + 24.5*x^2 + (343/6)*x^3 + (2401/24)*x^4 + (16807/120)*x^5 + (117649/720)*x^6) + (117649/720)*x^7 )};
            \addlegendentry{$c=7$}
            
            % c=8
            \addplot[MyPurple, very thick, domain=0:1, samples=100]
            {((131072/315)*x^8) / ( (1-x)*(1 + 8*x + 32*x^2 + (256/3)*x^3 + (512/3)*x^4 + (4096/15)*x^5 + (16384/45)*x^6 + (131072/315)*x^7) + (131072/315)*x^8 )};
            \addlegendentry{$c=8$}
            
        \end{axis}
    \end{tikzpicture}
    
    \vspace{2mm}
    {\small \textit{Зависимость вероятности ожидания от коэффициента загрузки $\rho$ при различном числе каналов. Все кривые асимптотически стремятся к $1$ при $\rho \to 1$.}}
\end{center}

\subsection*{Распределение времени ожидания и пребывания в системе M|M|c}

Аналогично системе с одним каналом, знание стационарного распределения и вероятности задержки позволяет получить функции распределения времени ожидания $W_q(t)$ и полного времени пребывания в системе $W(t)$ при условии использования дисциплины обслуживания FCFS.

\begin{statement}[Распределение времени ожидания и пребывания в M|M|c]
Для многоканальной системы M|M|c с дисциплиной обслуживания FCFS функции распределения времени ожидания $W_q(t) = \P(T_q \le t)$ и времени пребывания $W(t) = \P(T \le t)$ задаются выражениями (для $t \ge 0$):
\begin{itemize}
    \item $W_q(t) = 1 - C(c, r) e^{-c\mu(1-\rho)t}$;
    \item $W(t) = 1 - e^{-\mu t} - \dfrac{C(c, r)}{c(1-\rho) - 1} \left( e^{-\mu t} - e^{-c\mu(1-\rho)t} \right), \quad \text{если } c(1-\rho) \neq 1$.
\end{itemize}
В сингулярном случае, когда $c(1-\rho) = 1$, распределение полного времени имеет вид:
$$
W(t) = 1 - e^{-\mu t} - C(c, r) \mu t e^{-\mu t}.
$$
\end{statement}

\begin{sproof}
Пусть $N$ — число заявок в системе в момент прибытия новой заявки. Применим формулу полной вероятности:
$$
W_q(t) = \sum_{n=0}^{\infty} \P(T_q \le t \mid N = n) p_n.
$$
Если $n < c$, заявка немедленно поступает на свободный прибор, следовательно, $T_q = 0$. Сумма стационарных вероятностей для всех состояний до образования очереди равна $1 - \P(T_q > 0) = 1 - C(c, r)$.

Если $n \ge c$, заявка становится в очередь. Пусть $k = n - c \ge 0$ — число заявок, находящихся перед ней в очереди. Для того чтобы наша заявка перешла к обслуживанию, система должна завершить обработку $k+1$ заявок. Поскольку все $c$ каналов заняты, поток обслуживания является пуассоновским с интенсивностью $c\mu$. Время обработки $k+1$ заявок распределено по закону Эрланга порядка $k+1$ с параметром $c\mu$. Плотность этого распределения:
$$
f_{k+1}(x) = \frac{(c\mu)^{k+1}}{k!} x^k e^{-c\mu x}.
$$
Тогда функция распределения $W_q(t)$ принимает вид:
$$
W_q(t) = 1 - C(c, r) + \sum_{k=0}^{\infty} p_{c+k} \int_0^t \frac{(c\mu)^{k+1}}{k!} x^k e^{-c\mu x} dx.
$$
Используя свойство стационарных вероятностей $p_{c+k} = p_c \rho^k = C(c, r)(1-\rho)\rho^k$, поменяем порядок интегрирования и суммирования:
$$
W_q(t) = 1 - C(c, r) + \int_0^t C(c, r)(1-\rho) c\mu e^{-c\mu x} \left( \sum_{k=0}^{\infty} \frac{(c\mu\rho x)^k}{k!} \right) dx.
$$
Ряд сходится к экспоненте $e^{c\mu\rho x} = e^{\lambda x}$. Сводя показатели степеней, получаем интеграл:
$$
W_q(t) = 1 - C(c, r) + C(c, r) \int_0^t c\mu(1-\rho) e^{-c\mu(1-\rho)x} dx.
$$
Вычисляя интеграл, приходим к требуемому выражению:
$$
W_q(t) = 1 - C(c, r) + C(c, r) \left( 1 - e^{-c\mu(1-\rho)t} \right) = 1 - C(c, r) e^{-c\mu(1-\rho)t}.
$$

Для нахождения $W(t)$ воспользуемся тем, что полное время пребывания $T = T_q + T_{serv}$ есть сумма независимых случайных величин. Покажем, что плотность суммы независимых непрерывных величин $X$ и $Y$ вычисляется как свертка их плотностей. Вероятность попадания суммы в интервал определяется интегралом от совместной плотности по полуплоскости $x + y \le t$:
$$
F_{X+Y}(t) = \P(X+Y \le t) = \iint_{x+y \le t} f_{X,Y}(x,y) dx dy.
$$
В силу независимости $f_{X,Y}(x,y) = f_X(x)f_Y(y)$. Зафиксируем $x$ и проинтегрируем по $y$ от $-\infty$ до $t-x$:
$$
F_{X+Y}(t) = \int_{-\infty}^{\infty} f_X(x) \left( \int_{-\infty}^{t-x} f_Y(y) dy \right) dx = \int_{-\infty}^{\infty} f_X(x) F_Y(t-x) dx.
$$
Дифференцируя данное выражение по параметру $t$, получаем формулу свертки плотностей:
$$
f_{X+Y}(t) = \int_{-\infty}^{\infty} f_X(x) f_Y(t-x) dx = (f_X * f_Y)(t).
$$

Применим этот результат. Плотность времени обслуживания $f_{serv}(y) = \mu e^{-\mu y}$. Случайная величина $T_q$ принимает значение 0 с вероятностью $1 - C(c, r)$, поэтому ее плотность содержит дискретную компоненту. С использованием функции Дирака плотность времени ожидания записывается как:
$$
f_{T_q}(x) = (1 - C(c, r))\delta(x+0) + C(c, r) c\mu(1-\rho) e^{-c\mu(1-\rho)x}, \quad x \ge 0.
$$
Сдвиг аргумента дельта-функции $x+0$ означает, что центр распределения смещен бесконечно мало влево от нуля. Это необходимо для обеспечения непрерывности $W(t)$ в нуле справа. Для краткости обозначим $\gamma = c\mu(1-\rho)$. Подставим плотности в формулу свертки, интегрируя от $0-$ до $t$:
$$
f_T(t) = \int_{0-}^t \left[ (1 - C(c, r))\delta(x+0) + C(c, r) \gamma e^{-\gamma x} \right] \mu e^{-\mu(t-x)} dx.
$$
Интеграл распадается на два слагаемых. Первое вычисляется благодаря фильтрующему свойству функции Дирака, второе преобразуется путем вынесения экспоненты от $t$:
$$
f_T(t) = (1 - C(c, r))\mu e^{-\mu t} + C(c, r) \gamma \mu e^{-\mu t} \int_0^t e^{-(\gamma - \mu)x} dx.
$$

Рассмотрим общий случай $\gamma \neq \mu$ (то есть $c(1-\rho) \neq 1$). Вычисляя интеграл, получаем плотность:
$$
f_T(t) = (1 - C(c, r))\mu e^{-\mu t} + C(c, r) \frac{\gamma \mu}{\gamma - \mu} \left( e^{-\mu t} - e^{-\gamma t} \right).
$$
Интегрируя плотность $f_T(\tau)$ по отрезку $[0, t]$, находим функцию распределения полного времени $W(t)$:
$$
W(t) = \int_0^t f_T(\tau) d\tau = (1 - C(c, r))(1 - e^{-\mu t}) + C(c, r) \frac{\gamma \mu}{\gamma - \mu} \int_0^t \left( e^{-\mu \tau} - e^{-\gamma \tau} \right) d\tau.
$$
Беря первообразные, находим:
$$
W(t) = (1 - C(c, r))(1 - e^{-\mu t}) + C(c, r) \frac{\gamma}{\gamma - \mu} (1 - e^{-\mu t}) - C(c, r) \frac{\mu}{\gamma - \mu} (1 - e^{-\gamma t}).
$$
Сгруппируем слагаемые, содержащие множитель $(1 - e^{-\mu t})$:
$$
1 - C(c, r) + C(c, r) \frac{\gamma}{\gamma - \mu} = 1 + C(c, r) \left( \frac{\gamma}{\gamma - \mu} - 1 \right) = 1 + C(c, r) \frac{\mu}{\gamma - \mu}.
$$
Подставляя это обратно в выражение для $W(t)$ и раскрывая скобки:
\begin{align*}
    W(t) &= \left( 1 + C(c, r) \frac{\mu}{\gamma - \mu} \right) (1 - e^{-\mu t}) - C(c, r) \frac{\mu}{\gamma - \mu} (1 - e^{-\gamma t}) \\[2mm]
    &= 1 - e^{-\mu t} + C(c, r) \frac{\mu}{\gamma - \mu} - C(c, r) \frac{\mu}{\gamma - \mu} e^{-\mu t} - C(c, r) \frac{\mu}{\gamma - \mu} + C(c, r) \frac{\mu}{\gamma - \mu} e^{-\gamma t} \\[2mm]
    &= 1 - e^{-\mu t} - C(c, r) \frac{\mu}{\gamma - \mu} \left( e^{-\mu t} - e^{-\gamma t} \right).
\end{align*}
Подставляя обратно определение множителя $\frac{\mu}{\gamma - \mu} = \frac{\mu}{c\mu(1-\rho) - \mu} = \frac{1}{c(1-\rho) - 1}$, приходим к итоговой формуле для $W(t)$.

В сингулярном случае вырождения корней $\gamma = \mu$ интеграл в выражении для плотности вырождается в $\int_0^t 1 dx = t$. Плотность принимает вид распределения Эрланга:
$$
f_T(t) = (1 - C(c, r))\mu e^{-\mu t} + C(c, r) \mu^2 t e^{-\mu t}.
$$
Функция распределения находится интегрированием по частям $\left( \int \mu^2 \tau e^{-\mu \tau} d\tau = 1 - e^{-\mu t} - \mu t e^{-\mu t} \right)$:
\begin{align*}
    W(t) &= (1 - C(c, r))(1 - e^{-\mu t}) + C(c, r) \left( 1 - e^{-\mu t} - \mu t e^{-\mu t} \right) \\[2mm]
    &= (1 - e^{-\mu t}) - C(c, r) \mu t e^{-\mu t}.
\end{align*}
\end{sproof}

\begin{center}
    \pgfplotsset{compat=1.9}
    \begin{tikzpicture}
        \begin{axis}[
            width=0.9\textwidth, height=8cm,
            xlabel={Время $t$}, ylabel={$\P(\cdot \le t)$},
            xmin=0, xmax=8, ymin=0, ymax=1.05,
            grid=major,
            legend pos=south east,
            legend style={font=\tiny, cells={anchor=west}},
            legend columns=2,
            title={CDF времени ожидания ($W_q$) и пребывания ($W$) в $M|M|c$ ($\rho=0.75, \mu=1$)}
        ]
            % --- c=1 (DeepSkyBlue4) ---
            % W(t) = 1 - exp(-(1-0.75)*x) = 1 - exp(-0.25*x)
            \addplot[DeepSkyBlue4, thick, domain=0:8, samples=100] {1 - exp(-0.25*x)};
            \addlegendentry{$W, c=1$}
            % Wq(t) = 1 - 0.75*exp(-0.25*x)
            \addplot[DeepSkyBlue4, thick, dashed, domain=0:8, samples=100] {1 - 0.75*exp(-0.25*x)};
            \addlegendentry{$W_q, c=1$}

            % --- c=2 (Red3) ---
            % C(2, 1.5) = 0.6428. gamma = 2*1*0.25 = 0.5. 
            % W(t) = 1 - exp(-x) - (0.6428 / (2*0.25 - 1)) * (exp(-x) - exp(-0.5*x))
            % Т.к. 2(1-rho)-1 = -0.5, формула: 1 - exp(-x) + 1.2856*(exp(-x) - exp(-0.5*x))
            \addplot[Red3, thick, domain=0:8, samples=100] {1 - exp(-x) + 1.2856*(exp(-x) - exp(-0.5*x))};
            \addlegendentry{$W, c=2$}
            % Wq(t) = 1 - 0.6428*exp(-0.5*x)
            \addplot[Red3, thick, dashed, domain=0:8, samples=100] {1 - 0.6428*exp(-0.5*x)};
            \addlegendentry{$W_q, c=2$}

            % --- c=4 (SpringGreen4) ---
            % C(4, 3) = 0.5094. gamma = 4*1*0.25 = 1.0. 
            % ВНИМАНИЕ: Здесь gamma = mu = 1 (сингулярный случай).
            % W(t) = 1 - exp(-x) - 0.5094*x*exp(-x)
            \addplot[SpringGreen4, thick, domain=0:8, samples=100] {1 - exp(-x) - 0.5094*x*exp(-x)};
            \addlegendentry{$W, c=4$}
            % Wq(t) = 1 - 0.5094*exp(-1.0*x)
            \addplot[SpringGreen4, thick, dashed, domain=0:8, samples=100] {1 - 0.5094*exp(-x)};
            \addlegendentry{$W_q, c=4$}

            % Точки разрыва в нуле для Wq
            \draw[DeepSkyBlue4, fill=white] (axis cs:0, 0.25) circle (1.5pt); % 1-0.75
            \draw[Red3, fill=white] (axis cs:0, 0.3572) circle (1.5pt);       % 1-0.6428
            \draw[SpringGreen4, fill=white] (axis cs:0, 0.4906) circle (1.5pt); % 1-0.5094
            
        \end{axis}
    \end{tikzpicture}
\end{center}

\section{Оптимальный выбор числа каналов и режимы масштабирования}

Одной из ключевых задач при проектировании систем массового обслуживания является выбор числа обслуживающих приборов $c$ для заданной нагрузки $r = \lambda/\mu$. Этот выбор определяется компромиссом между качеством обслуживания (минимизацией вероятности ожидания $C(c, r)$) и эффективностью использования ресурсов (максимизацией коэффициента загрузки $\rho$).

Условие стационарности $c > r$ позволяет представить число каналов в виде $c = r + \Delta$, где $\Delta > 0$ — избыточный ресурс (резерв) производительности.

\begin{example}{Задача масштабирования}
Рассмотрим систему $M|M|c$ с нагрузкой $r_0 = 9$ Эрланг и числом каналов $c_0 = 12$. В исходном состоянии коэффициент загрузки $\rho_0 = 0.75$, а вероятность задержки $C(12, 9) \approx 0.266$.

Пусть нагрузка возрастает в 4 раза: $r_{new} = 36$. Необходимо определить новое число каналов $c_{new}$ исходя из различных стратегий управления ресурсами.
\end{example}

\begin{eproof}
Рассмотрим три асимптотических режима масштабирования при $r \to \infty$:

\begin{enumerate}
    \item \textbf{Режим качества (Quality Domain, QD).} 
    Фиксируется коэффициент загрузки: $\rho = \text{const} < 1$. Тогда $c_{QD} = r/\rho$. Для $r=36$ и $\rho=0.75$ получаем $c = 48$. При $r \to \infty$ вероятность задержки экспоненциально убывает $C(c, r) \to 0$. Система обладает избыточной надежностью, что экономически неэффективно из-за низкого использования ресурсов.

    \item \textbf{Режим эффективности (Efficiency Domain, ED).}
    Фиксируется абсолютный резерв каналов: $\Delta = c - r = \text{const}$. Тогда $c_{ED} = r + \Delta$. Для $r=36$ и $\Delta = 3$ получаем $c = 39$. При $r \to \infty$ загрузка $\rho \to 1$, однако вероятность задержки $C(c, r) \to 1$, что ведет к деградации качества обслуживания.

    \item \textbf{Сбалансированный режим (Quality-and-Efficiency Domain, QED).}
    Фиксируется вероятность обслуживания без задержек $C(c, r) = \text{const}$. Применение режима QED позволяет сохранить вероятность обслуживания без задержек на заданном уровне при одновременном повышении эффективности использования каналов. Известный как режим Халфина-Уитта, данный подход предполагает, что резерв каналов растет пропорционально квадратному корню из нагрузки:
    $$
    c_{QED} \approx r + \beta\sqrt{r}, \quad \beta > 0.
    $$
    Для исходной системы ($r_0=9, c_0=12$) параметр $\beta = (12 - 9)/\sqrt{9} = 1$. Тогда для новой нагрузки $r_{new}=36$ получаем:
    $$
    c_{new} \approx 36 + 1 \cdot \sqrt{36} = 42.
    $$
    При таком выборе вероятность задержки остается практически неизменной ($C(42, 36) \approx 0.27$), а коэффициент загрузки возрастает до $\rho = 36/42 \approx 0.857$.
\end{enumerate}
\end{eproof}

Математическое обоснование режима QED дает следующая предельная теорема.

\begin{theorem}{Теорема Халфина-Уитта (Halfin-Whitt)}
Рассмотрим последовательность систем $M|M|n$ с интенсивностью нагрузки $r_n$, где $n \to \infty$ и $r_n \to \infty$ так, что существует конечный предел:
$$
\lim_{n \to \infty} \frac{n - r_n}{\sqrt{r_n}} = \beta, \quad \beta > 0.
$$
Тогда предельная вероятность ожидания заявки в очереди существует и определяется соотношением:
$$
\lim_{n \to \infty} C(n, r_n) = \alpha(\beta) = \left[ 1 + \frac{\beta \Phi(\beta)}{\phi(\beta)} \right]^{-1},
$$
где $\phi(x)$ и $\Phi(x)$ — плотность и функция распределения стандартного нормального закона $\mathcal{N}(0,1)$ соответственно:
$$
\phi(x) = \frac{1}{\sqrt{2\pi}} e^{-\frac{x^2}{2}}, \quad \Phi(x) = \int_{-\infty}^{x} \phi(t) dt.
$$
\end{theorem}

\begin{note}
С практической точки зрения сбалансированный подход (QED) реализуется через решение задачи дискретной оптимизации. Если задан максимально допустимый уровень вероятности задержки в очереди (Quality of Service) $\alpha \in (0, 1)$, оптимальное число каналов определяется как:
$$
c^* = \min \{ c \in \N : C(c, r) \le \alpha \}.
$$
В режиме масштабирования Халфина-Уитта эта задача сводится к поиску соответствующего параметра $\beta$, который определяет необходимый избыток производительности (safety staffing).
\end{note}